\documentclass[../DoAn.tex]{subfiles}
\begin{document}

\section{Đặt vấn đề}
\label{section:1.1}

Ngành công nghiệp trò chơi điện tử trên toàn cầu đang trải qua giai đoạn tăng trưởng liên tục. Theo báo cáo thường niên của Newzoo \cite{Newzoo2024}, doanh thu thị trường game thế giới năm 2024 đạt khoảng 187,7 tỷ USD và dự kiến tiếp tục tăng trong các năm tới. Trong đó, phân khúc game trên nền tảng PC chiếm tỷ trọng đáng kể với khoảng 37,4 tỷ USD \cite{Newzoo2024}, phản ánh nhu cầu lớn của người dùng đối với các sản phẩm giải trí tương tác chất lượng cao trên nền tảng này. Trước đó, vào năm 2020, ngành game toàn cầu đã đạt mức doanh thu kỷ lục 179 tỷ USD, tăng 20\% so với năm 2019, cho thấy sức bền và khả năng tăng trưởng vượt trội so với nhiều ngành giải trí khác \cite{Witkowski2020}.

Trong số các thể loại trò chơi phổ biến, thể loại bắn súng hành động (Action Shooter) luôn duy trì vị trí hàng đầu về lượng người chơi và doanh thu. Các tựa game tiêu biểu như Control (Metascore 85/100 trên PC \cite{MetacriticControl}), Hitman: World of Assassination (Metascore 87/100 trên PC \cite{MetacriticHitman3}), và Max Payne (Metascore 89/100 trên PC \cite{MetacriticMaxPayne}) đã chứng minh sức hấp dẫn bền vững của thể loại bắn súng góc nhìn thứ ba (Third-Person Shooter -- TPS) qua nhiều thế hệ người chơi. Đặc biệt, những tựa game tích hợp các cơ chế sáng tạo như làm chậm thời gian (Bullet Time) trong Max Payne, hay sử dụng năng lực siêu nhiên để thao túng vật thể và không gian trong Control, đã nhận được phản hồi tích cực từ cả cộng đồng người chơi lẫn giới phê bình trên các trang đánh giá uy tín như Metacritic, bởi chúng mang lại trải nghiệm chiến thuật phong phú hơn so với các trò chơi bắn súng truyền thống. Tuy nhiên, trên thị trường hiện tại, số lượng sản phẩm kết hợp đồng thời các cơ chế bắn súng ẩn nấp, dịch chuyển không gian và thao túng thời gian trong một trò chơi TPS vẫn còn rất hạn chế.

Song song với nhu cầu của người chơi, sự phát triển nhanh chóng của thị trường game tại khu vực Đông Nam Á nói chung và Việt Nam nói riêng cũng mở ra triển vọng cho các nhà phát triển nội địa. Theo báo cáo của Niko Partners \cite{NikoPartners2024}, thị trường game Đông Nam Á năm 2024 có khoảng 270 triệu người chơi, trong đó Việt Nam đứng trong nhóm dẫn đầu về lượng người chơi game trên PC. Mặc dù phần lớn sản phẩm game phổ biến tại Việt Nam đến từ các nhà phát triển nước ngoài, nhu cầu và năng lực phát triển game nội địa đang dần được cải thiện. Việc một sinh viên hoặc nhà phát triển độc lập trong nước có thể tạo ra một sản phẩm game 3D hoàn chỉnh sẽ góp phần khẳng định tiềm năng của nguồn nhân lực công nghệ Việt Nam trong lĩnh vực này.

Bên cạnh đó, sự phát triển của các công cụ phát triển game thế hệ mới cũng mở ra cơ hội lớn cho các nhà phát triển độc lập (indie developer). Unreal Engine 5 với các công nghệ đồ họa tiên tiến như Nanite (Virtualized Geometry) và Lumen (Global Illumination) cho phép xây dựng môi trường 3D chất lượng cao mà không yêu cầu đội ngũ phát triển quy mô lớn. Đồng thời, hệ sinh thái tài nguyên mở ngày càng phong phú từ các nền tảng như Fab cung cấp hàng nghìn mô hình 3D, vật liệu và hoạt ảnh miễn phí, giúp giảm đáng kể chi phí sản xuất. Điều này đặt ra một bài toán thực tiễn: liệu một cá nhân có thể tận dụng các công cụ và tài nguyên hiện đại để phát triển một trò chơi hành động hoàn chỉnh với chi phí tối thiểu, đồng thời vẫn đảm bảo chất lượng trải nghiệm cho người chơi hay không.

Ngoài ra, từ góc độ học thuật, việc xây dựng một trò chơi điện tử hoàn chỉnh đòi hỏi sinh viên phải vận dụng tổng hợp nhiều lĩnh vực kiến thức: từ lập trình hệ thống, thiết kế kiến trúc phần mềm, xử lý toán học trong không gian 3D, đến thiết kế trải nghiệm người dùng và trí tuệ nhân tạo. Đây là một bài toán tổng hợp có giá trị thực tiễn cao, phù hợp với mục tiêu đào tạo của chương trình Khoa học máy tính, đồng thời cung cấp kinh nghiệm thực hành trực tiếp với một game engine mạnh mẽ và phổ biến trong ngành công nghiệp như Unreal Engine 5.

\section{Mục tiêu và phạm vi đề tài}
\label{section:1.2}

Trên thị trường game hành động bắn súng hiện tại, các sản phẩm phổ biến thường tập trung vào một trong hai hướng chính. Hướng thứ nhất là các tựa game bắn súng góc nhìn thứ ba với cơ chế ẩn nấp (Cover Shooter) như Gears of War (Metascore 94/100 \cite{MetacriticGearsOfWar}) hay The Division, nơi người chơi dựa vào hệ thống vật cản để chiến đấu. Hướng thứ hai là các tựa game khai thác cơ chế thao túng thời gian như Max Payne (Bullet Time, Metascore 89/100 \cite{MetacriticMaxPayne}) hay Superhot (thời gian chỉ trôi khi người chơi di chuyển). Một số sản phẩm khác như Portal và Splitgate (Metascore 71/100 \cite{MetacriticSplitgate}) lại khai thác cơ chế dịch chuyển không gian thông qua cổng không gian (Portal). Tuy nhiên, mỗi sản phẩm thường chỉ khai thác một cơ chế đơn lẻ, dẫn đến trải nghiệm chiến thuật bị giới hạn trong một khuôn khổ nhất định.

Về phía nhu cầu người chơi, cộng đồng game ngày càng tìm kiếm những trải nghiệm mới lạ, nơi nhiều cơ chế gameplay được kết hợp hài hòa để tạo ra chiều sâu chiến thuật. Theo báo cáo \textit{Essential Facts} của Entertainment Software Association (ESA) năm 2024 \cite{ESA2024}, khoảng 65\% người chơi cho biết họ ưu tiên lựa chọn trò chơi có cơ chế chơi sáng tạo và khác biệt so với các sản phẩm cùng thể loại. Điều này cho thấy tiềm năng của một sản phẩm có khả năng tích hợp đa cơ chế trong cùng một trải nghiệm.

Từ thực trạng trên, đồ án này đặt ra mục tiêu phát triển trò chơi Paradox Caliber -- một trò chơi bắn súng hành động góc nhìn thứ ba, trong đó người chơi có thể đồng thời sử dụng ba cơ chế cốt lõi: (i) bắn súng ẩn nấp (Cover Shooting) để chiến đấu dựa vào vật cản trong môi trường, (ii) cổng không gian (Portal) để dịch chuyển tức thời nhằm tạo lợi thế chiến thuật hoặc đánh úp đối phương, và (iii) làm chậm thời gian (Bullet Time) để né tránh đạn và phản ứng trong các tình huống nguy hiểm. Trò chơi được thiết kế với 07 màn chơi tuyến tính có bối cảnh và độ khó tăng dần, kèm theo hệ thống trí tuệ nhân tạo (AI) điều khiển kẻ thù có khả năng tự tìm chỗ nấp và phối hợp tấn công.

Lõi đồ án tận dụng tối đa các tài nguyên 3D miễn phí từ nền tảng Fab và các nguồn mở khác, kết hợp với các logic tinh chỉnh trong Unreal Engine 5.6 nhằm giảm thiểu tối đa chi phí sản xuất mà vẫn đạt chất lượng trải nghiệm cho đồ họa vượt trội.

Paradox Caliber hướng đến nhóm người chơi nam và nữ trong độ tuổi từ 16 đến 30, những người thường chơi các tựa game hành động bắn súng (TPS, FPS) trên nền tảng PC. Sản phẩm cũng thu hút nhóm yêu thích đề tài khoa học viễn tưởng (Sci-Fi) với các yếu tố như thao túng thời gian và du hành không gian. Về mặt sản phẩm, Paradox Caliber là trò chơi một người chơi (Single-player) trên PC. Mục tiêu hiệu năng là cung cấp đồ họa sắc nét trên khung hình mượt mà cho cấu hình tầm trung.

\section{Định hướng giải pháp}
\label{section:1.3}

Về thể loại, Paradox Caliber thuộc thể loại Third-Person Shooter (TPS) -- bắn súng hành động góc nhìn thứ ba, kết hợp các yếu tố Cover Shooter (bắn súng ẩn nấp). Đây là thể loại đã được kiểm chứng về mặt thị trường và có nền tảng cơ chế phù hợp để tích hợp thêm các yếu tố sáng tạo.

Về công nghệ, đồ án sử dụng Unreal Engine 5.6 làm nền tảng phát triển chính. Lý do lựa chọn được dựa trên: (i) hệ thống chiếu sáng động Lumen và Nanite, (ii) hệ thống AI Behavior Tree tích hợp sẵn làm cốt lõi trí tuệ nhân tạo, và (iii) hệ sinh thái tài nguyên mở của nhà phát triển Epic.

Giải pháp cụ thể của đồ án là xây dựng một trò chơi hoàn chỉnh với ba cơ chế hoạt động đồng thời. Cơ chế hệ thống bắn súng ẩn nấp quen thuộc; Cơ chế hệ thống cổng không gian (Portal), sử dụng các phép tính toán Vector và Rotator để thực hiện dịch chuyển tức thời; Và cơ chế hạ nhiệt nhịp độ thời gian (Bullet Time), sử dụng thuật toán Time Dilation.

\section{Bố cục đồ án}
\label{section:1.4}
Phần còn lại của báo cáo đồ án tốt nghiệp này được tổ chức như sau.

Chương 2 trình bày quá trình khảo sát và phân tích yêu cầu cho trò chơi Paradox Caliber. Nội dung chương bao gồm việc khảo sát các sản phẩm game hành động bắn súng tương tự trên thị trường để tìm ra các chức năng và mục đích chính của trò chơi.

Trong Chương 3, đồ án đi sâu vào chi tiết thiết kế trò chơi Paradox Caliber, bao gồm tổng quan lối chơi, các cơ chế cốt lõi như hệ thống điều khiển, cốt truyện cũng như hệ thống 5 màn chơi bối cảnh. Chương này kết hợp cùng phân loại và yêu cầu tương quan từ cơ chế không gian và thời gian.

Chương 4 cung cấp chi tiết về nền tảng lý thuyết và công nghệ sử dụng, giải thích sâu các nguyên lý vật lý, chiếu sáng trong Unreal Engine và cấu trúc cây hành vi AI Behavior Tree điều hướng kẻ địch.

Chương 5 và 6 trình bày kết quả thực nghiệm, giải pháp hệ thống và đánh giá về toàn bộ hoạt động thiết kế trò chơi, qua đó kết luận tiềm năng và định hướng phát triển ở tương lai.

\end{document}